\section{Running Example}\label{planta:sec:running_example}

Throughout this paper, the SUAVE exemplar~\cite{suave} is used as a running example to illustrate the concepts introduced by \acronym. In SUAVE, an Autonomous Underwater Vehicle (AUV) carries out an underwater pipeline inspection mission. The mission is structured as a sequence of task actions: \emph{A1}~search for the pipeline, and \emph{A2}~follow and inspect the pipeline. The SUAVE scenario is extended in ROSA~\cite{rosa} with a third action, \emph{A3}~recharge battery, so that the scenario requires task and architecture co-adaptation.

During mission execution, the AUV is subject to three sources of uncertainty. \Uone corresponds to thruster failures: if a thruster fails during \emph{A1} or \emph{A2}, the system requires reconfiguring the architecture to recover the affected thruster. \Utwo corresponds to the variation in water visibility: reduced visibility during \emph{A1} causes the system to select a lower-altitude search configuration. Both \Uone and \Utwo require architecture-level adaptation only, leaving the task plan unchanged. \Uthree corresponds to a critical drop in battery level: if battery level reaches a critical threshold during \emph{A1} or \emph{A2}, the system must interrupt the current action and execute \emph{A3}, which requires a dedicated architectural configuration. Hereafter, the scenario with \Uone and \Utwo is called the original use case, and the scenario with all three uncertainties the extended use case. % chktex 12 -- sentence starts with a macro

In this work, each task action in the mission is encoded as a PDDL action in the planning domain. For example, the \texttt{search\_pipeline} action~(Listing~
\ref{planta:lst:search_pipeline}) encodes \emph{A1}. Throughout the paper, this action is used to illustrate how \acronym transforms a regular action into a managed action that integrates functional and quality requirements derived from the \pddltomasys ontology~(Listing~\ref{planta:lst:search_pipeline_managed}), as described in Section~\ref{planta:sec:plan_transformation_taca}.
